\chapter{Implementation and Evaluation}
\label{chap:impl}

The preceding chapter defined the GEO pipeline at the architectural level.
This chapter documents how that architecture was translated into a working system,
following the execution order of the pipeline: from prompt generation through
multi-model querying, entity resolution, and web presence research, to the
experimental configuration and results. For each agent, the focus is on the
implementation decisions that the design phase left open and the non-trivial
engineering problems that arose during development.

% ═══════════════════════════════════════════════════════════════════════════════
\section{Development Environment and Technology Stack}
\label{sec:devenv}
% ═══════════════════════════════════════════════════════════════════════════════

The system was implemented in \textbf{Python 3.11} on a standard Windows 10 machine
without GPU hardware. This constraint directly shaped every technology choice:
all LLM inference had to be accessible via cloud APIs, all components had to operate
within free or low-cost tiers, and the tool integration layer had to remain
replaceable without touching agent logic.

Table~\ref{tab:stack} summarises the technologies selected under these constraints.
The three LLM providers were chosen to be complementary: Groq for speed,
OpenRouter for breadth of fallback models, and Mistral as a paid independent pool
for reliable structured output. The custom \texttt{ModelRegistry} ties them together
by managing quota state across all providers and keys autonomously.

\begin{table}[H]
\centering
\caption{Technology Stack}
\label{tab:stack}
\begin{tabular}{|p{3.2cm}|p{3.0cm}|p{6.8cm}|}
\hline
\textbf{Component} & \textbf{Technology} & \textbf{Role and Rationale} \\
\hline
Language & Python 3.11 & Native async support, dominant LLM tooling ecosystem \\
\hline
Orchestration & Custom ReAct + LangGraph & Adaptive execution, context-sensitive pipeline control \\
\hline
LLM Provider 1 & Groq API & Sub-second inference via GroqChip hardware; first-priority \\
\hline
LLM Provider 2 & OpenRouter & Unified access to 1,610+ models; primary fallback pool \\
\hline
LLM Provider 3 & Mistral API & Reliable JSON output; independent paid quota pool \\
\hline
Model Management & Custom \texttt{ModelRegistry} & Per-model quota tracking, autonomous provider rotation \\
\hline
Tool Protocol & FastMCP (stdio) & MCP-compliant tool servers, decoupled from agent code \\
\hline
Web Search & DuckDuckGo, SerpAPI & Complementary sources for entity web presence \\
\hline
Data Storage & pandas (CSV) & Incremental checkpointing, fault-tolerant resumption \\
\hline
Fuzzy Matching & RapidFuzz & C-accelerated Levenshtein distance for entity deduplication \\
\hline
Encoding & UTF-8 BOM (\texttt{utf-8-sig}) & French/Arabic character support on Windows CP1252 \\
\hline
\end{tabular}
\end{table}

% ═══════════════════════════════════════════════════════════════════════════════
\section{Agent 0: Prompt Generation and Self-Validation}
\label{sec:agent0impl}
% ═══════════════════════════════════════════════════════════════════════════════

Agent~0 generates the prompt set that drives the entire downstream analysis.
Its primary implementation challenge is quality control: an LLM left unconstrained
will produce prompts that are geographically vague, phrased to elicit generic
cuisine descriptions rather than specific brand names, or written in English
regardless of the requested language.

Three constraints are enforced via the system prompt. First, every prompt must be
geographically anchored to a specific Tunisian city — prompts drifting to Paris or
other non-Tunisian locations are structurally invalid for this study. Second, at
least 60\% of prompts must use formulations that force the responding LLM to recall
and name specific establishments — superlatives, rankings, and \textit{``cite-moi''}
constructions qualify; open-ended questions do not. Third, prompt text must be
written in the target language rather than defaulting to English.

These constraints are enforced not by post-hoc filtering but by a \textbf{self-reflection
loop}: after generating each intent's variants, Agent~0 evaluates its own output
against the quality criteria and returns a structured report including a quality score,
a brand-eliciting percentage, and a \texttt{proceed} flag. If the flag is false, the
agent regenerates with additional intent categories injected into the next iteration.
The loop runs for at most two cycles, after which the best available output is accepted.

A final guard ensures that an existing prompt CSV is never overwritten by a new
generation run. This protects a manually curated prompt set from being silently
replaced by a lower-quality LLM-generated one during pipeline restarts.

% ═══════════════════════════════════════════════════════════════════════════════
\section{Agent 1: Multi-Model Querying and Entity Processing}
\label{sec:agent1impl}
% ═══════════════════════════════════════════════════════════════════════════════

Agent~1 is the most computationally intensive stage, responsible for 450 LLM calls,
named entity extraction, entity deduplication, and GEO feature computation.
Its implementation raises three non-trivial decisions: how to preserve measurement
consistency across provider switches, how to resolve surface-form variation in
extracted entities, and how to quantify entity visibility in a way that captures
both frequency and positional stability.

\subsection{Slot Abstraction for Cross-Provider Consistency}

A core requirement of the GEO analysis is that entity visibility be measured
comparably across model capability tiers. If the pipeline records which specific
model and provider answered each query, a provider switch mid-run introduces a
confounding variable that conflates provider effects with model scale effects.

This is resolved through a \textbf{slot abstraction}: each model tier is assigned
a fixed label (\texttt{llama-small}, \texttt{qwen-medium}, \texttt{llama-large}),
and this label — not the actual model identifier — is what gets recorded in the output.
When Groq's daily quota is exhausted for \texttt{llama-large}, the registry
substitutes the OpenRouter equivalent at the same capability tier, but the slot
label remains unchanged. Per-tier GEO comparisons therefore remain valid regardless
of how many provider switches occurred during the run.

LLM parameters are also differentiated by role: querying runs at \texttt{temperature=0.7}
to encourage lexical diversity across the five runs; entity extraction runs at
\texttt{temperature=0.0} for deterministic, reproducible NER output.

\subsection{Two-Phase Entity Resolution}

Raw responses contain entities in highly variable surface forms: \textit{``Dar El Jeld''},
\textit{``dar el jeld''}, and \textit{``Restaurant Dar El Jeld''} must resolve to
the same canonical entity. Aggressive merging risks conflating distinct establishments;
under-merging fragments mention statistics across variants of the same entity.

The pipeline applies a two-phase resolution strategy.
\textbf{Phase B} uses \texttt{RapidFuzz}'s \texttt{token\_sort\_ratio}, which handles
word-order variance (\textit{``el ksar''} vs \textit{``ksar el''}) without penalising
legitimate transpositions. A score above 88 triggers automatic merge; a score between
75 and 87 flags the pair as ambiguous.

\textbf{Phase C} submits ambiguous pairs to an LLM arbitrator with a conservative
default of DIFFERENT — two names are merged only when the difference reduces to a
pure spelling variant, article prepend, or transliteration. The same call validates
entity legitimacy, rejecting establishments that are not Tunisian or North African
food-service businesses. This two-phase design keeps the fast local computation in
Phase B while delegating edge cases that require semantic reasoning to Phase C.

\subsection{GEO Visibility Scoring}

The central GEO metric is the \textbf{mention rate}: the fraction of all responses
in which an entity appears. Mention rate alone is insufficient, however — an entity
mentioned once at a consistent top position and one mentioned erratically across
many responses may have identical rates but very different GEO value.

The \textbf{stability score} addresses this by jointly weighting mention frequency
and ranking consistency:

\begin{equation}
\text{stability} \;=\; \text{mention\_rate} \;\times\; \frac{1}{1 + \sigma_{\text{rank}}}
\label{eq:stability}
\end{equation}

where $\sigma_{\text{rank}}$ is the standard deviation of ranking positions across
all mentions. Entities are classified as \textit{stable} (score $\geq 0.75$),
\textit{variable} ($0.40 \leq$ score $< 0.75$), or \textit{unstable} (score $< 0.40$).
Three complementary features complete the GEO profile of each entity:
\texttt{mention\_prominence} (fraction of mentions in top-3 positions),
\texttt{cross\_model\_rate} (fraction of model tiers that mention the entity), and
\texttt{co\_mention\_rate} (co-occurrence frequency with other top-ranked entities).

% ═══════════════════════════════════════════════════════════════════════════════
\section{Agent 2: Autonomous Web Presence Research}
\label{sec:agent2impl}
% ═══════════════════════════════════════════════════════════════════════════════

Agent~2 collects web presence data for each canonical entity by executing a
structured research protocol across seven sources: Google Maps, TripAdvisor,
Wikidata, official website, phone directory, Instagram, and Facebook.
The key implementation challenge is that this process cannot be scripted as a
fixed API call sequence — source availability varies per entity, results require
interpretation, and cascading dependencies exist between steps (a website must be
found before it can be scraped for social handles). A ReAct agent driving MCP tool
calls is therefore the appropriate mechanism.

\subsection{Tool-Capable Model Guard}

Agent~2 requires a model that reliably supports structured tool calling in a
multi-step ReAct loop. The \texttt{ModelRegistry} is designed to return any available
model as a fallback when the preferred model is quota-exhausted — a behaviour
appropriate for Agent~1 but dangerous for Agent~2, which would silently degrade
to a model incapable of tool calling.

A \texttt{\_TOOL\_CAPABLE} whitelist addresses this: when the registry returns a
model outside the list, it is rejected and the next candidate in the list is tried.
This guard is applied at both the initial model selection and at every model switch
triggered by quota exhaustion during processing, ensuring that Agent~2 never operates
on an untested model regardless of registry fallback behaviour.

\subsection{Confidence Scoring and Fault Tolerance}

At the end of each entity's research cycle, the agent produces an
\texttt{overall\_confidence} score according to explicit calibration guidelines:
0.0 indicates no usable data found; 0.5 indicates partial single-source coverage;
0.9 or above indicates multi-source verification including review data and at least
two independent signals. Delegating this assessment to the agent — rather than
computing it post-hoc from a fixed formula — allows the score to reflect data
consistency and plausibility, not just source count. A separate
\texttt{data\_source\_count} field, computed deterministically as a sum of binary
presence indicators, provides an objective cross-check.

Fault tolerance is enforced through a \texttt{done\_keys} filter at startup: only
rows with \texttt{confidence > 0} and no error are considered complete. Rows written
during quota failures are excluded and retried automatically on the next run.
The \texttt{SupervisorAgent} handles in-session failures by classifying each error
as either transient (TPM: wait and retry) or terminal (TPD: mark exhausted and
switch model). Table~\ref{tab:supervisor} details the classification rules.

\begin{table}[H]
\centering
\caption{Supervisor Error Classification Rules}
\label{tab:supervisor}
\begin{tabular}{|p{4.5cm}|p{3.0cm}|p{5.5cm}|}
\hline
\textbf{Error Signal} & \textbf{Action} & \textbf{Notes} \\
\hline
HTTP 429 + ``per day'' / ``free-models-per-day'' / HTTP 402 & \texttt{switch\_model} & Mark TPD-exhausted, persist timestamp to JSON \\
\hline
HTTP 429 + ``per minute'' / ``RPM'' & \texttt{wait\_retry} & Extract wait time from error message \\
\hline
HTTP 413 + token limit & \texttt{switch\_model} & Registry learns capacity, skips model for oversized prompts \\
\hline
HTTP 500 / 502 / 503 & \texttt{retry} & Transient server error, immediate retry \\
\hline
5+ consecutive RPM errors & \texttt{switch\_model} & Escape RPM loop; model effectively unusable \\
\hline
attempt $\geq 4$ & \texttt{skip} & Write error row with confidence~=~0, continue \\
\hline
\end{tabular}
\end{table}

% ═══════════════════════════════════════════════════════════════════════════════
\section{Engineering Challenges and Solutions}
\label{sec:challenges}
% ═══════════════════════════════════════════════════════════════════════════════

Beyond the agent-level decisions above, development surfaced a set of
system-wide failures that were only discoverable under live API conditions.
Table~\ref{tab:challenges} documents each failure, its root cause, and the fix applied.

\begin{table}[H]
\centering
\caption{Engineering Challenges and Implemented Solutions}
\label{tab:challenges}
\begin{tabular}{|p{3.6cm}|p{4.6cm}|p{4.6cm}|}
\hline
\textbf{Challenge} & \textbf{Root Cause} & \textbf{Solution} \\
\hline
Pipeline stalls on TPD exhaustion & No persistent per-model daily quota state & Registry persists exhaustion timestamps to JSON; survives restarts \\
\hline
Infinite RPM retry loop & Attempt counter decremented on wait-retry path, keeping it perpetually at 1 & Removed decrement; each wait consumes the attempt budget \\
\hline
\texttt{free-models-per-day} misclassified as TPM & Hyphenated string not matched by ``per day'' substring & Exact string added to TPD list; TPD evaluated before TPM \\
\hline
First CSV column name corrupted on Windows & Missing \texttt{utf-8-sig} codec prepends BOM bytes to first column & Consistent \texttt{utf-8-sig} on all read and write operations \\
\hline
French prompt set silently overwritten & LLM defaults to English on regeneration; no file guard & \texttt{agent0\_run()} never overwrites an existing prompt CSV \\
\hline
Error rows permanently skipped on resume & \texttt{\_load\_done()} counted all rows including failed ones & Filter: \texttt{confidence > 0} AND \texttt{error} is null \\
\hline
Agent 2 using non-tool-capable models & Registry returns any available fallback when preferred exhausted & Guard rejects models not in \texttt{\_TOOL\_CAPABLE} whitelist \\
\hline
8B models reject Agent 2 system prompt & System prompt exceeds small model's token budget & Registry learns capacity from 413 errors; skips model for oversized prompts \\
\hline
\end{tabular}
\end{table}

A common thread across these failures is that each one assumes a well-behaved external
environment: consistent error string formats, reliable codec defaults, and predictable
fallback behaviour. None of these assumptions held in practice, and none of the failures
were detectable through static analysis or unit testing alone.

% ═══════════════════════════════════════════════════════════════════════════════
\section{Experimental Setup}
\label{sec:setup}
% ═══════════════════════════════════════════════════════════════════════════════

\subsection{Domain and Motivation}

The experiment targets \textbf{Tunisian restaurants} across all cities and price
segments. This domain was selected for three reasons: it exhibits a large and directly
observable GEO visibility gap between internationally documented establishments and
those unknown outside Tunisia; it provides a realistic test of cross-lingual entity
normalisation due to the French-Arabic bilingual character of Tunisian online content;
and it extends GEO research to a North African, Francophone context that the
existing literature has not addressed~\cite{aggarwal2023geo}.

\subsection{Prompt Set and Query Configuration}

The study uses 30 French-language prompts across six intent categories
(Table~\ref{tab:prompts}), each evaluated under three model slots for five
independent runs:

\[
N \;=\; P \times S \times R \;=\; 30 \times 3 \times 5 \;=\; 450 \;\text{LLM queries}
\]

The three-tier slot design (\texttt{llama-small} 8B, \texttt{qwen-medium} 32B,
\texttt{llama-large} 70B) enables detection of scale-dependent visibility differences.
The five-run repetition supports stability estimation: an entity appearing consistently
across all runs of a given (prompt, slot) pair has high GEO reliability; one appearing
in a single run does not.

\begin{table}[H]
\centering
\caption{Prompt Set: Intent Categories and Examples}
\label{tab:prompts}
\begin{tabular}{|p{3.8cm}|p{4.5cm}|p{5.4cm}|}
\hline
\textbf{Intent} & \textbf{Description} & \textbf{Example Prompt} \\
\hline
\texttt{top\_recommendation} & Best restaurants by reputation & \textit{``Quel est le meilleur restaurant tunisien à Tunis ?''} \\
\hline
\texttt{cuisine\_authentique} & Traditional dish recommendations & \textit{``Quel restaurant est réputé pour son couscous à Tunis ?''} \\
\hline
\texttt{localisation} & Restaurants in a specific district & \textit{``Quels restaurants recommandes-tu dans le quartier La Marsa ?''} \\
\hline
\texttt{occasion\_speciale} & Context-specific venue recommendations & \textit{``Quel restaurant pour un dîner romantique à Tunis ?''} \\
\hline
\texttt{comparatif} & Explicit ranking and comparison & \textit{``Classe les 5 meilleurs restaurants tunisiens à Tunis.''} \\
\hline
\texttt{avis\_touristes} & Tourist-oriented queries & \textit{``Je suis touriste à Tunis, quel restaurant me recommandes-tu ?''} \\
\hline
\end{tabular}
\end{table}

% ═══════════════════════════════════════════════════════════════════════════════
\section{Results}
\label{sec:results}
% ═══════════════════════════════════════════════════════════════════════════════

\textit{Note: Results will be reported in full following pipeline completion.
The subsections below will be updated with per-agent quantitative analysis.}

\subsection{Agent 0 — Prompt Set Validation}
The self-reflection mechanism confirmed 100\% brand-elicitation validity across
all 30 prompts, six balanced intent categories, and coverage of seven Tunisian cities.

\subsection{Agent 1 — Response Collection and Entity Extraction}
\textit{[To be completed.]}

\subsection{Agent 2 — Web Presence Research}
\textit{[To be completed.]}

\subsection{Agent 3 — GEO Scoring and Visibility Ranking}
\textit{[To be completed.]}

% ═══════════════════════════════════════════════════════════════════════════════
\section{Discussion}
\label{sec:discussion}
% ═══════════════════════════════════════════════════════════════════════════════

\subsection{Architecture Validation}

Two architectural claims from Chapter~\ref{chap:design} received direct empirical
validation during the partial run. When Groq daily quota exhausted within minutes
of the first session, the registry substituted OpenRouter equivalents transparently,
with no impact on slot label consistency and no manual intervention — confirming
that the slot abstraction performs correctly under real provider switching conditions.
The ReAct orchestrator detected that a pre-built French prompt set already existed
and bypassed Agent~0 generation entirely, advancing directly to entity extraction.
This is precisely the kind of context-sensitive shortcut that fixed-graph pipelines
cannot perform.

\subsection{Limitations}

Three limitations bear on the current results.
Entity rankings from a partial query run are directional rather than definitive;
full 450-query coverage is required for confidence-interval-based comparisons.
The cleaning pipeline does not yet apply geographic validation, allowing
non-Tunisian establishments to persist as false positives — a fix that is
straightforward given the Wikidata \texttt{country} field already collected by
Agent~2.
Finally, five runs per (prompt, slot) pair produce substantial variance for
infrequent entities; ten or more runs would be preferable for tail-entity analysis.

% ═══════════════════════════════════════════════════════════════════════════════
\section{Chapter Summary}
% ═══════════════════════════════════════════════════════════════════════════════

This chapter followed the pipeline agent by agent, documenting the implementation
decisions that shaped system correctness and robustness.
Agent~0's self-reflection loop enforces prompt quality without manual curation.
Agent~1's slot abstraction decouples measurement identity from provider identity,
and its two-phase entity resolution balances speed with precision.
Agent~2's tool-capable model guard and supervisor-driven error classification
ensure that web research degrades gracefully under quota pressure rather than
silently failing.
The engineering challenges in Table~\ref{tab:challenges} illustrate a broader
point: resilient multi-provider LLM pipelines cannot be fully validated by design
alone — they require empirical testing under live quota and encoding conditions.
\end{document}
